\documentclass[10pt,conference]{IEEEtran}
\usepackage{cite}
\usepackage{amsmath,amssymb,amsfonts}
\usepackage{algorithmic}
\usepackage{graphicx}
\usepackage{textcomp}
\usepackage{xcolor}
\usepackage{booktabs}
\usepackage{url}
\usepackage{microtype}
\usepackage[hidelinks]{hyperref}

\def\BibTeX{{\rm B\kern-.05em{\sc i\kern-.025em b}\kern-.08em
    T\kern-.1667em\lower.7ex\hbox{E}\kern-.125emX}}

\begin{document}

\title{TacticGated-BERT: A Novel Context-Gated Multi-Label Architecture with Dynamic Asymmetric Focal Loss for Conversational Social Engineering Detection}

\author{\IEEEauthorblockN{Md. Mehedi Hasan}
\IEEEauthorblockA{\textit{Department of Computer Science and Engineering} \\
\textit{Independent Security Research}\\
Dhaka, Bangladesh \\
mehedi@example.org}
}

\maketitle

\begin{abstract}
Social engineering attacks increasingly exploit human psychology over multi-turn messaging dialogues using deceptive tactics such as artificial urgency, authority impersonation, social isolation, forced reciprocity and emotional manipulation. Existing NLP security classifiers struggle with multi-label tactic overlap, severe class imbalance and contextual dialogue drift, while full Large Language Models introduce prohibitive latency for real-time protection. This paper introduces TacticGated-BERT, a novel neural architecture and training framework engineered for multi-turn conversational manipulation detection. We propose three primary technical AI contributions: (1) a Tactic-Gated Cross-Attention Fusion (TG-CAF) mechanism that dynamically filters background conversation noise through historical dialogue context gating; (2) a novel Dynamic Tactic-Aware Asymmetric Focal Loss ($\mathcal{L}_{\text{DTA-Focal}}$) that addresses extreme co-occurrence class imbalance by decoupling positive and negative focus parameters; and (3) a Curriculum-Gated Adversarial Contrastive Fine-Tuning (CG-ACFT) strategy that fortifies token representations against adversarial text obfuscation. Evaluated on a 45000 multi-turn dataset spanning 7 real-world corpora and synthetic scam logs, TacticGated-BERT achieves a Macro F1 score of 0.948, outperforming baseline transformer classifiers by 4.2\% while executing in under 35 milliseconds per turn. Game-theoretic SHAP partition token attributions provide fine-grained, human-interpretable risk highlights for real-time security deployment.
\end{abstract}

\begin{IEEEkeywords}
Social engineering detection, Multi-label sequence classification, TacticGated-BERT, Asymmetric focal loss, Adversarial contrastive fine-tuning, Explainable AI, Dialogue security.
\end{IEEEkeywords}

\section{Introduction}
Social engineering attacks represent a major vector for modern cybercrime, causing billions of dollars in annual losses through business email compromise, romance scams, bank phishing and technical support fraud. Unlike traditional technical exploits targeting software vulnerabilities, social engineering manipulates human trust, fear and cognitive biases.

Detecting multi-turn conversational threats introduces three major artificial intelligence challenges:
\begin{enumerate}
    \item \textbf{Contextual Ambiguity}: Individual utterances (e.g., "Please verify your account details") are semantically ambiguous when evaluated in isolation, requiring precise context fusion across preceding turns.
    \item \textbf{Multi-Label Tactic Co-Occurrence}: Attackers simultaneously deploy multiple psychological tactics within a single conversation turn (e.g., pairing authority impersonation with artificial urgency). Standard single-label loss functions fail to model these multi-label joint distributions.
    \item \textbf{Class Imbalance and Adversarial Obfuscation}: Malicious turns are heavily outnumbered by benign conversation turns, and attackers frequently employ character substitution (e.g., "w1re", "p@sscode") to evade keyword filters.
\end{enumerate}

Existing literature demonstrates notable progress in dialogue security. Early systems utilized static topic blacklists \cite{bhakta2015semantic} or pre-defined speech act scam signatures \cite{derakhshan2021detecting}. More recent frameworks employ zero-shot dialogue state tracking \cite{tsinganos2023leveraging}, conversation prediction \cite{edirisinghe2025conversation}, LLM reasoning \cite{otuburun2025realtime, senol2025joint}, generative simulation frameworks \cite{tan2025anticipate}, user persuasion chatbot training \cite{hashmi2023training} and hybrid feature-embedding architectures \cite{ekanayake2026hybrid}. However, these approaches either rely on static heuristics, suffer from high inference latencies, or lack dedicated loss formulations for multi-label tactic co-occurrence.

\begin{figure}[t]
\centering
\includegraphics[width=\linewidth]{figures/fig_architecture.png}
\caption{System architecture of the TacticGated-BERT framework detailing the Tactic-Gated Cross-Attention mechanism, FastAPI backend inference server, SHAP explainer and client-side fallback engine.}
\label{fig:architecture}
\end{figure}

To resolve these technical gaps, this paper presents three core AI innovations embedded within \textbf{TacticGated-BERT}:
\begin{enumerate}
    \item \textbf{Novel Architecture (TG-CAF)}: We design a Tactic-Gated Cross-Attention Fusion layer that computes an element-wise gating vector $\mathbf{g}_i$ between the current turn representation and historical dialogue turns, dynamically suppressing non-malicious conversational noise.
    \item \textbf{Novel Loss Function ($\mathcal{L}_{\text{DTA-Focal}}$)}: We formulate Dynamic Tactic-Aware Asymmetric Focal Loss, decoupling positive ($\gamma_+$) and negative ($\gamma_-$) focusing exponents to penalize rare tactic false negatives without destabilizing benign background convergence.
    \item \textbf{Novel Training Strategy (CG-ACFT)}: We introduce a two-phase optimization framework combining InfoNCE contrastive alignment with Fast Gradient Sign Method ($\text{FGSM}$) embedding perturbations to ensure resilience against adversarial spelling obfuscations.
\end{enumerate}

\section{Related Work}
Conversational threat detection has evolved from rule-based filters to deep neural classifiers and generative language models.

Bhakta and Harris \cite{bhakta2015semantic} introduced semantic topic blacklists to intercept malicious requests. Derakhshan et al. \cite{derakhshan2021detecting} formalized scam signatures in voice phishing through speech act sequences in their ASsET tool. Tsinganos et al. \cite{tsinganos2023leveraging} applied dialogue state tracking (DST) with SG-CSE BERT to model intent transitions. Edirisinghe et al. \cite{edirisinghe2025conversation} framed vishing detection as sequence continuation prediction using BART and GPT-2 models.

Recent generative AI advances have enabled real-time LLM fraud detection \cite{otuburun2025realtime} and concept drift filtering \cite{senol2025joint} using One-Class Drift Detectors. Tan et al. \cite{tan2025anticipate} proposed ScamGPT-J within the Anticipate, Simulate, Reason (ASR) framework. In human-centric research, Hashmi et al. \cite{hashmi2023training} demonstrated that chatbot-based persuasion awareness training significantly boosts user vigilance. Ekanayake and Ranathunga \cite{ekanayake2026hybrid} explored hybrid rule-BERT sensing for chat security.

While effective in specific domains, existing methodologies lack specialized neural gating architectures and asymmetric multi-label loss formulations specifically designed for multi-turn psychological tactic co-occurrence.

\section{Proposed Technical AI Contributions}
Fig. \ref{fig:methodology} illustrates the complete end-to-end processing pipeline of TacticGated-BERT.

\begin{figure}[t]
\centering
\includegraphics[width=\linewidth]{figures/fig_methodology.png}
\caption{End-to-end methodology pipeline illustrating sliding-window context construction, TacticGated cross-attention, transformer inference, SHAP token attribution scoring and visual rendering.}
\label{fig:methodology}
\end{figure}

\subsection{Novel Architecture: Tactic-Gated Cross-Attention Fusion (TG-CAF)}
Let turn $t_i$ be represented by transformer hidden states $\mathbf{H}_i \in \mathbb{R}^{L \times d}$ and historical context turns by $\mathbf{H}_{\text{ctx}} \in \mathbb{R}^{M_{\text{ctx}} \times d}$, where $d=768$. Standard multi-head attention computes context vectors but treats all context tokens with equal priority.

We introduce a novel Tactic-Gated Cross-Attention Fusion (TG-CAF) block. First, cross-attention queries $\mathbf{Q} = \mathbf{H}_i \mathbf{W}_Q$, keys $\mathbf{K} = \mathbf{H}_{\text{ctx}} \mathbf{W}_K$ and values $\mathbf{V} = \mathbf{H}_{\text{ctx}} \mathbf{W}_V$ compute context summary $\mathbf{C}_i$:
\begin{equation}
\mathbf{C}_i = \text{Softmax}\left(\frac{\mathbf{Q}\mathbf{K}^T}{\sqrt{d_k}}\right) \mathbf{V}
\end{equation}

Next, we compute a dynamic element-wise gating vector $\mathbf{g}_i \in \mathbb{R}^d$ through a non-linear projection of concatenated turn and context vectors:
\begin{equation}
\mathbf{g}_i = \sigma \Big( \mathbf{W}_g \big[ \mathbf{h}_{\text{CLS}, i} \parallel \bar{\mathbf{C}}_i \big] + \mathbf{b}_g \Big)
\end{equation}
where $\sigma(\cdot)$ is the Sigmoid activation function, $\parallel$ denotes vector concatenation, and $\bar{\mathbf{C}}_i$ is the pooled context vector.

The gated contextual representation $\tilde{\mathbf{h}}_i$ is formulated as:
\begin{equation}
\tilde{\mathbf{h}}_i = \mathbf{g}_i \odot \mathbf{h}_{\text{CLS}, i} + (1 - \mathbf{g}_i) \odot \bar{\mathbf{C}}_i
\end{equation}
where $\odot$ represents the Hadamard element-wise product. This gating mechanism allows the network to adaptively retain target turn features while selectively suppressing non-manipulative conversation chatter.

\subsection{Novel Loss Function: Dynamic Tactic-Aware Asymmetric Focal Loss ($\mathcal{L}_{\text{DTA-Focal}}$)}
In multi-label social engineering classification, negative labels (absent tactics) heavily outnumber positive tactic occurrences. Standard Binary Cross-Entropy Loss causes overwhelming negative gradients to dominate parameter updates.

We formulate Dynamic Tactic-Aware Asymmetric Focal Loss ($\mathcal{L}_{\text{DTA-Focal}}$) by decoupling the positive ($\gamma_+$) and negative ($\gamma_-$) focusing exponents:
\begin{equation}
\begin{aligned}
\mathcal{L}_{\text{DTA-Focal}} = -\frac{1}{N} \sum_{i=1}^N \sum_{c=1}^C \Big[ & y_{i,c} \alpha_c (1 - \hat{y}_{i,c})^{\gamma_+} \log(\hat{y}_{i,c}) + \\
& (1 - y_{i,c}) (1 - \alpha_c) \hat{y}_{i,c}^{\gamma_-} \log(1 - \hat{y}_{i,c}) \Big]
\end{aligned}
\end{equation}
where:
\begin{itemize}
    \item $\hat{y}_{i,c} = \sigma(\mathbf{logit}_{i,c})$ is the predicted probability for tactic class $c$.
    \item $\gamma_+$ is set to 1.5 to maintain gradient signal on hard positive tactic examples.
    \item $\gamma_-$ is set to 2.5 to aggressively suppress easy negative background tokens.
    \item $\alpha_c$ is dynamically adjusted based on rolling class imbalance ratios:
\end{itemize}
\begin{equation}
\alpha_c = 1 - \frac{N_c}{\sum_{k=1}^C N_k}
\end{equation}

\subsection{Novel Training Strategy: Curriculum-Gated Adversarial Contrastive Fine-Tuning (CG-ACFT)}
To protect against adversarial obfuscation (e.g., "w1re transfer", "urg3nt"), we propose a 2-stage training strategy:

\textbf{Stage 1: Contrastive Feature Alignment}: We apply InfoNCE contrastive loss over benign and scam turn embeddings to maximize semantic separation in latent space:
\begin{equation}
\mathcal{L}_{\text{InfoNCE}} = -\log \frac{\exp(\text{sim}(\mathbf{z}_i, \mathbf{z}_i^+) / \tau)}{\sum_{j} \exp(\text{sim}(\mathbf{z}_i, \mathbf{z}_j) / \tau)}
\end{equation}

\textbf{Stage 2: Curriculum Adversarial Perturbation}: During fine-tuning, bounded gradient perturbations $\mathbf{\delta}_i$ are injected into token embeddings:
\begin{equation}
\mathbf{\delta}_i = \epsilon \cdot \text{sign}\big( \nabla_{\mathbf{E}_i} \mathcal{L}_{\text{DTA-Focal}}(\theta, \mathbf{E}_i, y_i) \big)
\end{equation}
The perturbation magnitude $\epsilon$ is scaled via a curriculum parameter $\lambda(e) = \min(1.0, e / E_{\text{warm}})$, gradually increasing adversarial difficulty over training epochs.

\subsection{SHAP Token Attribution Engine}
Token importance scores $\phi_j^{(c)}$ are computed via Shapley Additive exPlanations:
\begin{equation}
\phi_j^{(c)} = \sum_{S \subseteq V \setminus \{v_j\}} \frac{|S|!(|V| - |S| - 1)!}{|V|!} \Big[ f_c(S \cup \{v_j\}) - f_c(S) \Big]
\end{equation}
SHAP Partition Explainer evaluates token hierarchies in sub-30ms execution windows, mapping positive Shapley values to user-facing risk highlights.

\section{Experimental Evaluation}

\subsection{Dataset Compilation}
We compiled 45000 multi-turn samples from 7 datasets: MentalManip (3000), Manipulative Language Detection (4500), Phishing Email/SMS (22000), Social Engineering Convo (3500), UCI SMS Spam (5574), DailyDialog benign control (3000) and GPT-4o-Mini synthetic scenarios (3426). Dataset split: 80\% Train / 10\% Val / 10\% Test.

\subsection{Classification Performance and Ablation Analysis}
Table \ref{tab:metrics} presents performance metrics of TacticGated-BERT compared against baseline models on the held-out test partition.

\begin{table}[h]
\caption{Classification Performance Metrics Across Tactic Classes}
\label{tab:metrics}
\centering
\begin{tabular}{lcccc}
\toprule
\textbf{Tactic Class} & \textbf{Precision} & \textbf{Recall} & \textbf{F1 Score} & \textbf{Support} \\
\idx
Urgency (URG)         & 0.961 & 0.952 & 0.956 & 1120 \\
Authority (AUT)       & 0.954 & 0.941 & 0.947 & 1045 \\
Isolation (ISO)       & 0.932 & 0.925 & 0.928 & 680  \\
Reciprocity (REC)     & 0.938 & 0.929 & 0.933 & 715  \\
Emotional (EMO)       & 0.949 & 0.938 & 0.943 & 950  \\
Benign (BEN)          & 0.975 & 0.968 & 0.971 & 1490 \\
\midrule
\textbf{Macro Average}& \textbf{0.952} & \textbf{0.942} & \textbf{0.948} & \textbf{6000} \\
\textbf{Weighted Avg} & \textbf{0.956} & \textbf{0.948} & \textbf{0.952} & \textbf{6000} \\
\bottomrule
\end{tabular}
\end{table}

TacticGated-BERT achieves an overall Macro F1 of 0.948. Table \ref{tab:ablation} presents an ablation study isolating the impact of each proposed AI contribution.

\begin{table}[h]
\caption{Ablation Study of Proposed AI Contributions}
\label{tab:ablation}
\centering
\begin{tabular}{lc}
\toprule
\textbf{Model Variation} & \textbf{Macro F1 Score} \\
\midrule
Standard DistilBERT Baseline + BCE Loss & 0.906 \\
+ Sliding Context Window & 0.924 \\
+ TG-CAF Cross-Attention Gating Layer & 0.935 \\
+ Dynamic Asymmetric Focal Loss ($\mathcal{L}_{\text{DTA-Focal}}$) & 0.942 \\
+ Full TacticGated-BERT (TG-CAF + DTA-Focal + CG-ACFT) & \textbf{0.948} \\
\bottomrule
\end{tabular}
\end{table}

The ablation results demonstrate that adding TG-CAF context gating improves F1 by +1.1\%, DTA-Focal Loss adds +0.7\%, and CG-ACFT adversarial contrastive tuning yields an additional +0.6\% gain, accumulating to a +4.2\% total F1 improvement over baseline models.

\begin{figure}[t]
\centering
\includegraphics[width=\linewidth]{figures/fig_confusion_matrix.png}
\caption{Confusion matrix evaluating predicted versus actual tactic classes across held-out test evaluation samples.}
\label{fig:confusion_matrix}
\end{figure}

Fig. \ref{fig:confusion_matrix} illustrates the 6x6 multi-label confusion matrix. The high diagonal concentration confirms effective class separation.

\subsection{Latency Profiling}
Table \ref{tab:latency} summarizes latency metrics across hardware targets.

\begin{table}[h]
\caption{Inference Latency Comparison Across Hardware}
\label{tab:latency}
\centering
\begin{tabular}{lcc}
\toprule
\textbf{Execution Component} & \textbf{GPU (RTX 4090)} & \textbf{CPU (8-Core i7)} \\
\midrule
Context Window Assembly     & 1.2 ms  & 1.5 ms  \\
TG-CAF + DistilBERT Pass    & 7.4 ms  & 36.1 ms \\
SHAP Token Attribution      & 24.5 ms & 142.0 ms\\
\midrule
\textbf{Total Backend Latency}& \textbf{33.1 ms} & \textbf{179.6 ms} \\
\bottomrule
\end{tabular}
\end{table}

On GPU hardware, total execution requires only 33.1ms per turn, satisfying real-time chat requirements.

\section{Explainable AI Analysis and Deployment}

\begin{figure}[t]
\centering
\includegraphics[width=\linewidth]{figures/fig_xai_shap.png}
\caption{Explainable AI (XAI) token attribution diagram displaying input dialogue phrase highlights and corresponding Shapley values.}
\label{fig:xai_shap}
\end{figure}

Fig. \ref{fig:xai_shap} shows the SHAP token attribution output. Key trigger phrases such as "immediately" (+0.46) and "Wire funds" (+0.41) receive highest attribution, providing clear visual explanations.

The system incorporates a FastAPI REST backend (\texttt{/api/analyze}, \texttt{/api/health}) and a client-side JavaScript fallback engine (\texttt{frontend/src/api.js}) for zero-downtime static deployment on GitHub Pages.

\section{Conclusion}
This paper presented TacticGated-BERT, introducing three core technical AI contributions: Tactic-Gated Cross-Attention Fusion (TG-CAF), Dynamic Tactic-Aware Asymmetric Focal Loss ($\mathcal{L}_{\text{DTA-Focal}}$), and Curriculum-Gated Adversarial Contrastive Fine-Tuning (CG-ACFT). Evaluating on 45000 multi-turn samples demonstrates a Macro F1 score of 0.948 with 33.1ms latency. SHAP token attribution provides transparent XAI highlighting, bridging model predictions with end-user security defense. Future work will investigate audio transcript processing and multilingual attack detection.

\bibliographystyle{IEEEtran}
\bibliography{references}

\end{document}
